\documentclass[10pt,a4paper]{article}
\usepackage[margin=17mm,headheight=13pt,headsep=6mm,footskip=9mm]{geometry}
\usepackage{amsmath,amssymb,amsthm}
\usepackage[T1]{fontenc}
\usepackage{lmodern,microtype}
\usepackage{enumitem,fancyhdr}
\usepackage[hidelinks]{hyperref}
\hypersetup{pdftitle={D3 Architect: Range GCD},pdfauthor={Anushri Maheshwari (250041005)}}
\setlength{\parindent}{0pt}
\setlength{\parskip}{4pt}
\setlist[itemize]{leftmargin=*,nosep}
\pagestyle{fancy}
\fancyhf{}
\fancyhead[L]{\small D3 Architect: Range GCD}
\fancyhead[R]{\small Anushri Maheshwari \,|\, 250041005}
\fancyfoot[C]{\small\thepage\ of 2}
\renewcommand{\headrulewidth}{0.3pt}
\newcommand{\g}{\operatorname{gcd}}
\newcommand{\sect}[1]{\par\smallskip{\large\bfseries #1}\par\smallskip}
\begin{document}
\begin{center}
{\LARGE\bfseries Range GCD: Algebra and Query Complexity}\\[4pt]
{\large Anushri Maheshwari \quad Roll No.: 250041005}\\[2pt]
CS 253 / MA 253 \quad Project 5 \quad D3 Architect
\end{center}
\vspace{-3pt}
\textbf{Setting.}
Store nonnegative integers in an array of length $n$, with point assignments
and zero-based half-open queries $[l,r)$.
For signed inputs, take their nonnegative magnitudes before the gcd fold.
Define $\g(0,0)=0$.
Let $V$ bound every value throughout the replay, including later assignments.
Set $H=1+\log_2(n+1)$ and $L=1+\log_2(V+1)$ to cover empty arrays and zeros.

\sect{1. GCD is a commutative monoid}
For each positive integer $d$, the defining divisibility property is
\[
 d\mid \g(a,b)\quad\Longleftrightarrow\quad d\mid a\ \text{and}\ d\mid b.
\]
This includes zero inputs, since every positive integer divides zero.
Nonnegative integers with the same positive divisors are equal: positive
ones divide each other, while only zero has every positive integer as a divisor.

\textbf{Commutativity.}
The common-divisor condition is symmetric, so $\g(a,b)=\g(b,a)$.

\textbf{Associativity.}
For every positive $d$,
\[
\begin{aligned}
d\mid\g(\g(a,b),c)
&\Longleftrightarrow d\mid a,\ d\mid b,\ d\mid c\\
&\Longleftrightarrow d\mid\g(a,\g(b,c)).
\end{aligned}
\]
The characterization gives $\g(\g(a,b),c)=\g(a,\g(b,c))$, also when zeros occur.

\textbf{Identity.}
$\g(a,0)$ and $a$ have the same positive divisors, so
$\g(a,0)=a=\g(0,a)$. Thus $(\mathbb{Z}_{\ge0},\g,0)$ is a commutative monoid.

\sect{2. Segment-tree correctness and the product bound}
Each real leaf stores the current value, padding leaves store $0$, and
each parent stores the gcd of its children.
Induction on height proves that each node holds the gcd of its covered leaves.
An assignment recomputes its ancestor path bottom to top, restoring the invariant
after increases, decreases, or replacement by zero.

A range query selects disjoint node intervals partitioning $[l,r)$.
Associativity permits their gcds to be folded, and identity $0$ handles empty
ranges and padding. The implementation uses two ordered accumulators.
At most two blocks are selected at each of $O(H)$ levels; traversal also costs
$O(H)$. Euclid requires $O(L)$ remainder operations per gcd call.
Therefore the direct bound is
\[
 \boxed{O(HL)}
 \qquad\text{or }O(\log n\cdot\log V)\text{ when }n,V\ge2.
\]

\sect{3. A ratio-sensitive Euclidean lemma}
\textbf{Lemma.}
For positive $a,b$ and $d=\g(a,b)$, remainder-based Euclid performs
$O(1+\log_2(a/d))$ divisions, even if $b>a$.

\textbf{Proof.}
Divide both inputs and all nonzero remainders by $d$.
The quotients and number of divisions are unchanged.
After at most one initial division, the smaller positive operand is at most $a/d$.
For $x\ge y>0$, write $r=x\bmod y$.
If $r\le y/2$, the smaller operand has already halved.
Otherwise $r>y/2$, and the next remainder is $y\bmod r=y-r<y/2$.
Thus it halves within at most two divisions until termination.
This gives the claimed logarithmic bound.
If $b=0$, the call returns $a$ in constant time instead. \hfill$\square$

\newpage
\sect{4. The tighter bound by Euclidean telescoping}
First consider one accumulator receiving node gcds $x_1,\ldots,x_k$.
Put $g_0=0$ and $g_i=\g(g_{i-1},x_i)$.
While the accumulator is zero, a call costs $O(1)$:
$\g(0,x)$ takes at most one remainder operation.
If no positive value is encountered, the entire fold costs $O(k)$.

Otherwise let $g_j$ be the first positive accumulator; then $1\le g_j\le V$.
For $i>j$, $g_i$ is a positive divisor of $g_{i-1}$.
The lemma gives a cost of
\[
 O\!\left(1+\log_2\frac{g_{i-1}}{g_i}\right)
\]
for call $i$. The same expression covers $x_i=0$, which leaves the
accumulator unchanged at constant cost. Summing the logarithmic charges,
\[
\begin{aligned}
\sum_{i=j+1}^{k}\log_2\frac{g_{i-1}}{g_i}
 &=\log_2\left(
     \frac{g_j}{g_{j+1}}\frac{g_{j+1}}{g_{j+2}}\cdots
     \frac{g_{k-1}}{g_k}\right)\\
 &=\log_2\frac{g_j}{g_k}
 \le \log_2 V.
\end{aligned}
\]
The empty sum covers $j=k$.
Thus the fold takes $O(k+L)$ operations.
Counting only the number of strict accumulator decreases would be insufficient:
the lemma also pays for the divisions \emph{inside} each call.

\textbf{Match to the code.}
The left and right query accumulators each form such a chain.
Their calls may be interleaved, but they telescope independently.
There are $O(H)$ selected blocks in total, and the final merge of the
accumulators adds at most $O(L)$. Including traversal yields
\[
 \boxed{O(H+L)=O(1+\log(n+1)+\log(V+1))}.
\]
For $n,V\ge2$ this is conventionally written
$\boxed{O(\log n+\log V)}$.
This is an upper bound, not a claim that every query attains it.

\sect{5. Computational model and limits}
\begin{itemize}
\item \textbf{Arithmetic-operation model.}
The proof counts Euclidean remainder steps and tree operations.
It is a time bound on a word RAM only when values, indices, and addresses
fit machine words and comparison, indexing, and division/remainder have constant cost.
The word width must accommodate both $\log(V+1)$ and $\log(n+1)$.
\item \textbf{Fixed-width C++.}
The application restricts calibration factors to $[0,10^9]$ and stores them
as \texttt{unsigned long long}. For a fixed value bound, $L$ is constant
and queries cost $O(H)$ in this model.
Actual remainder latency depends on the machine.
Signed backend inputs are converted to unsigned magnitudes safely, including
the minimum signed value.
\item \textbf{Bit complexity.}
For arbitrary-precision $b$-bit inputs, with
$b=\max\{1,\lceil\log_2(V+1)\rceil\}$, there are $O(H+b)$ divisions;
they are not constant-time.
Charging a conservative $O(b^2)$ bit operations per schoolbook remainder
gives $O((H+b)b^2)$ arithmetic bit operations, plus representation and index costs.
This deliberately loose estimate is not an optimal bit bound.
\end{itemize}

\sect{6. Implementation scope and other bounds}
The gcd tree uses $O(n)$ storage and $O(n)$ construction merges, giving the
safe construction bound $O(nL)$.
Retain the safe $O(HL)$ point-assignment bound:
recomputing sibling pairs is not the same accumulator chain as a query.
The additive proof specifically concerns the explicit Euclidean gcd routine.

In \texttt{app/}, both replay modes obtain calibration gcd from Structure A.
The B-hybrid mode uses Structure B only for maximum and A for maximum count
and gcd; it is not a pure-B gcd implementation.
Differential and invariant tests check the code, while the argument above
establishes the asymptotic query bound.

\smallskip
{\small\textbf{References.}
Project 5 brief, D3 requirement.
D.\ E.\ Knuth, \emph{The Art of Computer Programming}, Vol.\ 2, 3rd ed.,
\S4.5.3 (Euclid's algorithm).
The ratio lemma and query telescope are proved explicitly above.}
\end{document}
